% Setting up the document and importing necessary packages
\documentclass[12pt,a4paper]{article}

% Chinese support packages
\usepackage{xeCJK}
\usepackage{fontspec}
\setCJKmainfont{Noto Serif CJK TC}
\setCJKsansfont{Noto Sans CJK TC}
\setCJKmonofont{Noto Sans CJK TC}

% Other necessary packages
\usepackage{geometry}
\usepackage{titlesec}
\usepackage{enumitem}
\usepackage{color}
\usepackage{colortbl}
\usepackage{tcolorbox}
\usepackage{booktabs}
\usepackage{longtable}
\usepackage{array}
\usepackage{graphicx}
\usepackage{tikz}
\usepackage{mdframed}
\usepackage{fancyhdr}
\usepackage{hyperref}
\usepackage{multicol}
\usepackage{datetime2}
\usepackage{natbib}

\hypersetup{
    colorlinks=true,
    linkcolor=rheum_purple,
    citecolor=rheum_purple,
    urlcolor=rheum_purple,
    pdfborder={0 0 0}
}

% Page setup
\geometry{
    top=2.5cm,
    bottom=2.5cm,
    left=2cm,
    right=2cm,
    headheight=15pt
}

% Color definitions (Rheumatology themed colors)
\definecolor{rheum_purple}{RGB}{138,43,226}
\definecolor{rheum_blue}{RGB}{75,0,130}
\definecolor{rheum_gray}{RGB}{120,120,120}
\definecolor{highlight_yellow}{RGB}{255,250,205}

% Title format settings
\titleformat{\section}[block]{\Large\bfseries\color{rheum_purple}}{\thesection}{10pt}{}
\titleformat{\subsection}[block]{\large\bfseries\color{rheum_blue}}{\thesubsection}{8pt}{}
\titleformat{\subsubsection}[block]{\normalsize\bfseries\color{rheum_gray}}{\thesubsubsection}{6pt}{}

% Custom environment
\newmdenv[
    linecolor=rheum_purple,
    backgroundcolor=highlight_yellow,
    linewidth=2pt,
    topline=true,
    bottomline=true,
    leftline=true,
    rightline=true,
    innertopmargin=10pt,
    innerbottommargin=10pt,
    innerleftmargin=10pt,
    innerrightmargin=10pt
]{importantbox}

\newcommand{\note}[1]{
    \par
    \noindent
    \begin{tcolorbox}[
        colback=rheum_blue!5!white,
        colframe=rheum_blue,
        title=\textbf{重點提醒},
        fonttitle=\bfseries,
        boxrule=0.5pt,
        arc=3pt,
        left=6pt,
        right=6pt,
        top=6pt,
        bottom=6pt
    ]
    #1
    \end{tcolorbox}
}

% Header and footer settings
\pagestyle{fancy}
\fancyhf{}
\fancyhead[L]{\textcolor{rheum_purple}{內科部住院醫師教學文件}}
\fancyhead[R]{\textcolor{rheum_blue}{風濕免疫科EPAs}}
\fancyfoot[C]{\textcolor{rheum_gray}{\thepage}}
\renewcommand{\headrulewidth}{1pt}
\renewcommand{\headrule}{\hbox to\headwidth{\color{rheum_purple}\leaders\hrule height \headrulewidth\hfill}}

% Date format
\DTMsetstyle{yyyy-mm-dd}
\newcommand{\mydate}{\DTMtoday}

% Document start
\begin{document}

% Title page
\begin{titlepage}
    \centering
    {\LARGE\textcolor{rheum_purple}{\textbf{內科部住院醫師教學文件}}\par}
    \vspace{1cm}
    {\huge\bfseries 風濕免疫科可信賴專業活動EPAs\par}
    \vspace{0.5cm}
    {\Large\textcolor{rheum_blue}{\textit{Rheumatology Entrustable Professional Activities}}\par}
    \vspace{1cm}
    
    \begin{tcolorbox}[
        colback=rheum_blue!5!white,
        colframe=rheum_blue,
        width=0.8\textwidth,
        arc=10pt,
        boxrule=1pt
    ]
    \centering
    {\large\textbf{目標對象}\par}
    \vspace{0.3cm}
    內科部住院醫師R1-R3\par
    風濕免疫科專科訓練醫師\par
    \vspace{0.5cm}
    {\large\textbf{文件版本}\par}
    \vspace{0.3cm}
    第1.0版(10個EPAs)\par
    發布日期:\mydate\par
    \end{tcolorbox}

    \vfill
    
    {\large 風濕免疫科  \\
    適用於CCC評量系統\par}
    
    \vspace{0.5cm}
\end{titlepage}

\newpage
\tableofcontents
\newpage

% Main content
\section{風濕免疫科EPAs概述}

\subsection{EPA定義}
Entrustable Professional Activities (EPA) 是指可信賴委託給住院醫師執行的專業活動單元。風濕免疫科EPAs涵蓋了風濕免疫疾病診斷、治療與照護的核心能力,是所有風濕免疫科醫師必須精熟的專業技能。

\vspace{0.5cm}
\note{
風濕免疫科EPAs評量著重於能否安全、穩定地執行完整的風濕免疫疾病診療流程,並能整合資訊做出適當臨床決策。每個EPA都包含明確的里程碑發展階段,從初學者到專家的完整能力建構。本版本包含10個核心EPA,從基礎的病史詢問與關節檢查,到複雜的免疫調節藥物管理與跨科部協調,提供完整的風濕免疫疾病診療能力架構。
}

\subsection{學習目標}
完成風濕免疫科EPAs後,住院醫師應能夠:
\begin{itemize}[nosep]
    \item 熟練進行風濕免疫疾病病史詢問與風險評估。
    \item 準確執行關節系統完整理學檢查。
    \item 判讀風濕免疫疾病相關影像檢查。
    \item 執行關節穿刺術與關節液分析。
    \item 獨立判讀自體免疫抗體與免疫學檢驗。
    \item 診斷與治療類風濕性關節炎。
    \item 診斷與管理全身性紅斑狼瘡與結締組織疾病。
    \item 診斷與處置痛風與結晶性關節炎。
    \item 使用與監測免疫調節藥物(DMARDs與生物製劑)。
    \item 協調風濕免疫疾病患者跨科部整合照護。
\end{itemize}

\subsection{委託等級}
\begin{table}[h]
    \centering
    \caption{EPA委託等級定義}
    \begin{tabular}{p{2cm} p{3cm} p{9cm}}
        \toprule
        \textbf{等級} & \textbf{委託程度} & \textbf{描述} \\
        \midrule
        Level 1 & 觀察學習 & 僅能觀察他人執行,無法獨立操作 \\
        Level 2 & 直接監督 & 可在主治醫師直接監督下執行 \\
        Level 3 & 間接監督 & 可獨立執行,但需回報並接受指導 \\
        Level 4 & 監督可得 & 可獨立執行,監督者隨時可得 \\
        Level 5 & 完全獨立 & 可獨立執行並指導他人 \\
        \bottomrule
    \end{tabular}
\end{table}

\newpage
\section{風濕免疫科10大EPAs詳述}

\begin{longtable}{|p{1.2cm}|p{3.8cm}|p{8.5cm}|}
\caption{風濕免疫科EPAs完整架構} \\
\hline
\rowcolor{rheum_blue!20}
\textbf{EPA} & \textbf{專業活動名稱} & \textbf{核心能力要素與里程碑} \\
\hline
\endfirsthead

\multicolumn{3}{c}%
{{\bfseries \tablename\ \thetable{} -- 接續上頁}} \\
\hline
\rowcolor{rheum_blue!20}
\textbf{EPA} & \textbf{專業活動名稱} & \textbf{核心能力要素與里程碑} \\
\hline
\endhead

\hline \multicolumn{3}{|r|}{{接續下頁}} \\ \hline
\endfoot

\hline \hline
\endlastfoot

\rowcolor{rheum_purple!10}
\multicolumn{3}{|c|}{\textbf{EPA 1: 風濕免疫疾病病史詢問與風險評估}} \\
\hline
\textbf{描述} & \multicolumn{2}{|p{12.5cm}|}{能夠獨立進行風濕免疫疾病相關的完整病史詢問,包括症狀評估、家族史調查、藥物史,並能初步評估疾病活性與風險。} \\
\hline
\textbf{核心能力} & 症狀評估 & 關節痛、晨間僵硬、肌肉無力、皮疹、口腔潰瘍的詳細詢問 \\
\cline{2-3}
& 系統性症狀 & 發燒、體重減輕、疲倦、雷諾氏現象評估 \\
\cline{2-3}
& 家族史識別 & 自體免疫疾病、乾癬、僵直性脊椎炎家族史 \\
\cline{2-3}
& 藥物史 & NSAIDs、類固醇、免疫抑制劑使用史 \\
\cline{2-3}
& 功能評估 & 日常生活活動能力、工作能力評估(HAQ, RAPID3) \\
\cline{2-3}
& 疾病活性評估 & DAS28、SLEDAI等疾病活性指標應用 \\
\hline
\textbf{Level 1} & Novice & 能詢問基本關節症狀,需要監導者引導完成詢問 \\
\hline
\textbf{Level 2} & Advanced Beginner & 能系統性完成風濕病史詢問,使用適當的症狀評估工具 \\
\hline
\textbf{Level 3} & Competent & 能獨立完成複雜病史詢問,整合症狀與家族史進行初步風險分層 \\
\hline
\textbf{Level 4} & Proficient & 能在困難情況下獲取準確病史,整合多重共病資訊 \\
\hline
\textbf{Level 5} & Expert & 能指導他人進行病史詢問,識別罕見症狀表現 \\
\hline

\rowcolor{rheum_purple!10}
\multicolumn{3}{|c|}{\textbf{EPA 2: 關節系統理學檢查}} \\
\hline
\textbf{描述} & \multicolumn{2}{|p{12.5cm}|}{能夠熟練且準確地進行全身關節系統的完整理學檢查,包括關節腫脹、壓痛、活動度評估,並能正確解讀檢查發現。} \\
\hline
\textbf{核心能力} & 關節檢查 & 對稱性、腫脹、壓痛、溫度、積液檢查 \\
\cline{2-3}
& 活動度評估 & 主動與被動關節活動度、疼痛弧評估 \\
\cline{2-3}
& 特殊檢查 & 擠壓測試、骶髂關節、脊椎活動度評估 \\
\cline{2-3}
& 關節外表現 & 皮疹、結節、指甲病變、黏膜病變 \\
\cline{2-3}
& 關節計數 & 腫脹關節計數、壓痛關節計數(66/68 joints) \\
\cline{2-3}
& 功能測試 & 握力、行走速度、Schober test \\
\hline
\textbf{Level 1} & Novice & 能進行基本關節觸診,識別明顯腫脹 \\
\hline
\textbf{Level 2} & Advanced Beginner & 能識別常見異常理學發現,進行系統性關節檢查 \\
\hline
\textbf{Level 3} & Competent & 能區分各種異常理學發現,準確計數關節,整合多項檢查發現 \\
\hline
\textbf{Level 4} & Proficient & 能識別複雜理學發現特徵,檢查技術精準且有效率 \\
\hline
\textbf{Level 5} & Expert & 能偵測細微的異常發現,整合檢查發現與臨床表現 \\
\hline

\rowcolor{rheum_purple!10}
\multicolumn{3}{|c|}{\textbf{EPA 3: 風濕免疫疾病影像判讀}} \\
\hline
\textbf{描述} & \multicolumn{2}{|p{12.5cm}|}{能夠獨立判讀風濕免疫疾病相關影像檢查,包括關節X光、超音波、MRI,識別正常與異常變化,並提出臨床相關性分析。} \\
\hline
\textbf{核心能力} & 關節X光判讀 & 關節間隙狹窄、骨侵蝕、骨質增生、鈣化 \\
\cline{2-3}
& 超音波檢查 & 滑膜增生、血流訊號、積液、肌腱病變 \\
\cline{2-3}
& MRI判讀 & 骨髓水腫、滑膜炎、侵蝕、軟骨變化 \\
\cline{2-3}
& 分期評估 & Steinbrocker分期、Sharp score評估 \\
\cline{2-3}
& 鑑別診斷 & 發炎性vs退化性關節炎影像特徵 \\
\cline{2-3}
& 追蹤評估 & 疾病進展、治療反應影像評估 \\
\hline
\textbf{Level 1} & Novice & 能識別基本正常關節X光結構 \\
\hline
\textbf{Level 2} & Advanced Beginner & 能識別明顯異常病變,判讀基本超音波影像 \\
\hline
\textbf{Level 3} & Competent & 能獨立判讀大部分影像異常,整合影像與臨床症狀 \\
\hline
\textbf{Level 4} & Proficient & 能判讀複雜影像變化,識別細微的早期病變 \\
\hline
\textbf{Level 5} & Expert & 成為影像判讀的專家,能教導複雜影像判讀 \\
\hline

\rowcolor{rheum_purple!10}
\multicolumn{3}{|c|}{\textbf{EPA 4: 關節穿刺術與關節液分析}} \\
\hline
\textbf{描述} & \multicolumn{2}{|p{12.5cm}|}{能夠執行關節穿刺術,包括適應症評估、解剖定位、無菌技術,並能正確分析關節液檢查結果。} \\
\hline
\textbf{核心能力} & 適應症評估 & 診斷性vs治療性穿刺決策、禁忌症識別 \\
\cline{2-3}
& 解剖定位 & 膝、肩、踝、腕、肘等關節穿刺點定位 \\
\cline{2-3}
& 無菌技術 & 消毒、局部麻醉、穿刺技巧 \\
\cline{2-3}
& 關節液分析 & 外觀、白血球計數、結晶檢查、培養 \\
\cline{2-3}
& 結晶判讀 & 尿酸結晶、焦磷酸鈣結晶偏光顯微鏡判讀 \\
\cline{2-3}
& 併發症處理 & 出血、感染、疼痛管理 \\
\hline
\textbf{Level 1} & Novice & 能協助關節穿刺,了解基本解剖結構 \\
\hline
\textbf{Level 2} & Advanced Beginner & 能在監督下執行大關節穿刺 \\
\hline
\textbf{Level 3} & Competent & 能獨立執行關節穿刺,判讀關節液結果 \\
\hline
\textbf{Level 4} & Proficient & 能執行困難關節穿刺,使用超音波導引 \\
\hline
\textbf{Level 5} & Expert & 成為關節穿刺專家,能教導複雜技術 \\
\hline

\rowcolor{rheum_purple!10}
\multicolumn{3}{|c|}{\textbf{EPA 5: 自體免疫抗體與免疫學檢驗判讀}} \\
\hline
\textbf{描述} & \multicolumn{2}{|p{12.5cm}|}{能夠獨立判讀自體免疫抗體與免疫學檢驗報告,包括ANA、ENA、RF、ACPA等,識別正常與異常變化,並提出臨床相關性分析。} \\
\hline
\textbf{核心能力} & ANA判讀 & 效價、型態(speckled, homogeneous, nucleolar)判讀 \\
\cline{2-3}
& 特異性抗體 & Anti-dsDNA, Anti-Sm, Anti-RNP, Anti-Ro/La \\
\cline{2-3}
& RF與ACPA & 類風濕因子、抗環瓜氨酸抗體臨床意義 \\
\cline{2-3}
& 補體系統 & C3、C4、CH50在SLE的應用 \\
\cline{2-3}
& 發炎指標 & ESR、CRP在疾病活性評估 \\
\cline{2-3}
& 臨床關聯 & 抗體與疾病表現、預後的關聯性 \\
\hline
\textbf{Level 1} & Novice & 能識別明顯陽性抗體,知道基本正常範圍 \\
\hline
\textbf{Level 2} & Advanced Beginner & 能判讀常見自體抗體,使用系統性判讀流程 \\
\hline
\textbf{Level 3} & Competent & 能獨立判讀複雜抗體組合,整合臨床表現 \\
\hline
\textbf{Level 4} & Proficient & 能判讀罕見抗體型態,準確評估臨床意義 \\
\hline
\textbf{Level 5} & Expert & 成為免疫檢驗判讀專家,能教導複雜判讀 \\
\hline

\rowcolor{rheum_purple!10}
\multicolumn{3}{|c|}{\textbf{EPA 6: 類風濕性關節炎診斷與治療}} \\
\hline
\textbf{描述} & \multicolumn{2}{|p{12.5cm}|}{能夠診斷類風濕性關節炎,評估疾病活性與嚴重度,制定個別化治療計畫,包括DMARDs與生物製劑選擇,並監測治療效果。} \\
\hline
\textbf{核心能力} & 診斷標準 & 2010 ACR/EULAR分類標準應用 \\
\cline{2-3}
& 疾病活性評估 & DAS28、SDAI、CDAI評分系統 \\
\cline{2-3}
& 預後評估 & 侵蝕性疾病、關節外表現識別 \\
\cline{2-3}
& 治療策略 & Treat-to-target策略、藥物選擇階梯 \\
\cline{2-3}
& DMARDs使用 & MTX、HCQ、SSZ等使用與監測 \\
\cline{2-3}
& 併發症管理 & 關節變形、肺纖維化、血管炎處理 \\
\hline
\textbf{Level 1} & Novice & 能識別RA症狀徵象,知道基本診斷標準 \\
\hline
\textbf{Level 2} & Advanced Beginner & 能使用標準治療指引,進行基本藥物調整 \\
\hline
\textbf{Level 3} & Competent & 能獨立管理穩定RA,調整藥物治療 \\
\hline
\textbf{Level 4} & Proficient & 能處理複雜RA,評估進階治療選項 \\
\hline
\textbf{Level 5} & Expert & 成為RA專家,發展個人化治療策略 \\
\hline

\rowcolor{rheum_purple!10}
\multicolumn{3}{|c|}{\textbf{EPA 7: 全身性紅斑狼瘡與結締組織疾病管理}} \\
\hline
\textbf{描述} & \multicolumn{2}{|p{12.5cm}|}{能夠診斷全身性紅斑狼瘡與其他結締組織疾病,評估器官侵犯程度,制定個別化治療方案,並監測疾病活性與藥物副作用。} \\
\hline
\textbf{核心能力} & 診斷評估 & SLICC/ACR標準、EULAR/ACR 2019分類標準 \\
\cline{2-3}
& 器官侵犯評估 & 腎、肺、心、神經、血液系統評估 \\
\cline{2-3}
& 疾病活性 & SLEDAI、BILAG評分系統應用 \\
\cline{2-3}
& 免疫抑制治療 & 類固醇、MMF、CYC、AZA使用 \\
\cline{2-3}
& 生物製劑 & Belimumab、Rituximab適應症與使用 \\
\cline{2-3}
& 併發症預防 & 感染、骨質疏鬆、心血管疾病預防 \\
\hline
\textbf{Level 1} & Novice & 能識別SLE可能性,了解基本檢查 \\
\hline
\textbf{Level 2} & Advanced Beginner & 能執行基本評估,安排初步檢查 \\
\hline
\textbf{Level 3} & Competent & 能獨立診斷SLE,制定治療計畫 \\
\hline
\textbf{Level 4} & Proficient & 能處理複雜SLE,整合內外科治療 \\
\hline
\textbf{Level 5} & Expert & 成為SLE專家,處理罕見疾病與併發症 \\
\hline

\rowcolor{rheum_purple!10}
\multicolumn{3}{|c|}{\textbf{EPA 8: 痛風與結晶性關節炎診斷與處置}} \\
\hline
\textbf{描述} & \multicolumn{2}{|p{12.5cm}|}{能夠診斷痛風與其他結晶性關節炎,執行急性發作處置,制定長期降尿酸治療計畫,並監測治療效果與副作用。} \\
\hline
\textbf{核心能力} & 臨床診斷 & ACR/EULAR痛風分類標準、鑑別診斷 \\
\cline{2-3}
& 結晶確認 & 關節液偏光顯微鏡檢查、結晶判讀 \\
\cline{2-3}
& 急性發作處置 & NSAIDs、秋水仙素、類固醇選擇 \\
\cline{2-3}
& 降尿酸治療 & Allopurinol、Febuxostat使用與監測 \\
\cline{2-3}
& 共病管理 & 慢性腎病、高血壓、代謝症候群整合治療 \\
\cline{2-3}
& 預防策略 & 飲食建議、藥物預防、生活型態調整 \\
\hline
\textbf{Level 1} & Novice & 能識別典型痛風症狀,知道基本治療 \\
\hline
\textbf{Level 2} & Advanced Beginner & 能診斷與治療單純性痛風 \\
\hline
\textbf{Level 3} & Competent & 能獨立管理痛風,執行降尿酸治療 \\
\hline
\textbf{Level 4} & Proficient & 能處理難治性痛風,整合多重因素 \\
\hline
\textbf{Level 5} & Expert & 成為痛風專家,發展個人化治療策略 \\
\hline

\rowcolor{rheum_purple!10}
\multicolumn{3}{|c|}{\textbf{EPA 9: 免疫調節藥物使用與監測}} \\
\hline
\textbf{描述} & \multicolumn{2}{|p{12.5cm}|}{能夠合理選用免疫調節藥物,包括傳統DMARDs與生物製劑,了解其作用機轉、副作用,並執行適當監測與調整。} \\
\hline
\textbf{核心能力} & DMARDs選擇 & MTX、HCQ、SSZ、LEF藥理與適應症 \\
\cline{2-3}
& 生物製劑 & Anti-TNF、Anti-IL6、Anti-CD20、JAK抑制劑 \\
\cline{2-3}
& 副作用監測 & 肝腎功能、血球、感染風險監測 \\
\cline{2-3}
& 感染預防 & 疫苗接種、TB篩檢、帶狀皰疹預防 \\
\cline{2-3}
& 藥物交互作用 & 與其他藥物的交互作用評估 \\
\cline{2-3}
& 懷孕用藥 & 孕期安全藥物選擇與調整 \\
\hline
\textbf{Level 1} & Novice & 能了解常用免疫藥物,知道基本副作用 \\
\hline
\textbf{Level 2} & Advanced Beginner & 能使用基本DMARDs,執行標準監測 \\
\hline
\textbf{Level 3} & Competent & 能獨立使用免疫藥物,調整劑量與監測 \\
\hline
\textbf{Level 4} & Proficient & 能使用生物製劑,處理複雜副作用 \\
\hline
\textbf{Level 5} & Expert & 成為免疫藥物專家,制定個人化用藥策略 \\
\hline

\rowcolor{rheum_purple!10}
\multicolumn{3}{|c|}{\textbf{EPA 10: 風濕免疫疾病跨科部整合照護}} \\
\hline
\textbf{描述} & \multicolumn{2}{|p{12.5cm}|}{能夠協調風濕免疫疾病患者的整體照護,包括跨科部會診、多專科團隊合作、復健安排,並能有效溝通與教育病患及家屬。} \\
\hline
\textbf{核心能力} & 會診協調 & 腎臟科、皮膚科、復健科會診安排 \\
\cline{2-3}
& 團隊合作 & 多專科團隊會議參與、治療計畫整合 \\
\cline{2-3}
& 復健安排 & 物理治療、職能治療、輔具評估 \\
\cline{2-3}
& 病患教育 & 疾病衛教、藥物使用、自我監測指導 \\
\cline{2-3}
& 品質改善 & 照護品質監測、改善方案執行 \\
\cline{2-3}
& 家屬溝通 & 病情解釋、治療選項討論、預後說明 \\
\cline{2-3}
& 倫理考量 & 醫療倫理、安寧照護決策 \\
\hline
\textbf{Level 1} & Novice & 能參與基本病患照護,了解團隊角色 \\
\hline
\textbf{Level 2} & Advanced Beginner & 能執行基本會診,參與團隊討論 \\
\hline
\textbf{Level 3} & Competent & 能獨立協調病患照護,制定整合計畫 \\
\hline
\textbf{Level 4} & Proficient & 能整合複雜多科照護,領導團隊討論 \\
\hline
\textbf{Level 5} & Expert & 成為照護協調專家,改善照護品質 \\
\hline

\end{longtable}

\newpage
\section{評量方法與標準}

\subsection{CCC評量架構}
本風濕免疫科EPAs採用Case-based Clinical Competency (CCC)評量方式:
\begin{itemize}[nosep]
    \item 真實臨床情境評量
    \item 多次觀察與回饋
    \item 漸進式委託決策
    \item 360度回饋機制
\end{itemize}

\begin{importantbox}
\textbf{特別提醒:} 風濕免疫科EPAs的評量重點在於漸進式能力發展,而非單次完美表現。鼓勵從錯誤中學習,持續改進。每個EPA都需要在真實臨床情境中反復練習與評量。
\end{importantbox}

\subsection{評量工具對應表}
\begin{table}[h]
    \centering
    \caption{各EPA推薦評量工具}
    \begin{tabular}{p{4.5cm} p{4cm} p{5.5cm}}
        \toprule
        \textbf{EPA類型} & \textbf{主要評量工具} & \textbf{輔助評量方法} \\
        \midrule
        技能操作類 & DOPS & 直接觀察、技能檢核表 \\
        (EPA 4) & (Direct Observation of Procedural Skills) & 同儕評量、自我評量 \\
        \midrule
        臨床推理類 & Mini-CEX & 病例討論、口試 \\
        (EPA 1, 5, 6, 7, 8, 9) & (Mini-Clinical Evaluation Exercise) & 病歷撰寫評量、模擬情境 \\
        \midrule
        理學檢查類 & DOPS & 標準化病人、技能測驗 \\
        (EPA 2) & (Direct Observation of Procedural Skills) & 影片回顧、同儕觀察 \\
        \midrule
        影像判讀類 & 影像判讀評量 & 線上測驗、病例討論 \\
        (EPA 3) & Image Interpretation Assessment & 定期考核、peer review \\
        \midrule
        整合照護類 & Multi-source Feedback & 跨科團隊評量 \\
        (EPA 10) & (360-degree Assessment) & 病人滿意度調查 \\
        \bottomrule
    \end{tabular}
\end{table}

\section{總結}

風濕免疫科EPAs涵蓋了風濕免疫疾病診療的完整核心能力,從基礎的病史詢問與關節檢查,到複雜的免疫調節藥物管理、生物製劑使用,以及整合性的跨科部協調照護。每個EPA都有清楚的能力描述與里程碑發展路徑,協助住院醫師循序漸進地建立專業能力。

\begin{importantbox}
\textbf{關鍵訊息:} 風濕免疫科EPAs的核心在於「可委託性」與「完整性」。這10個EPA代表風濕免疫科醫師必須具備的基礎診療能力。關節檢查是風濕免疫科的基石,需要反復練習才能精熟。免疫調節藥物的使用需要謹慎評估與密切監測,確保病患安全。
\end{importantbox}

\end{document}